%% ============================================================
%%  SECCIÓN 6 — Procedimiento Experimental
%%  Responsable: Vera Vivanco, Jesús Rodolfo
%% ============================================================

\section{Procedimiento Experimental}

El procedimiento llevado a cabo para la determinación del coeficiente de dilatación lineal se detalla de forma cronológica a continuación:

\begin{enumerate}
    \item \textbf{Medición de la aguja de rodamiento:} Se utilizó el calibrador vernier para determinar con precisión el diámetro $D$ de la aguja metálica cilíndrica. El valor del radio $R$ se calculó como $D/2$, y se registró junto a su incertidumbre instrumental correspondiente en la hoja de recolección de datos.
    \item \textbf{Instalación de la muestra en el dilatómetro:} Se colocó el tubo de ensayo (iniciando con el cobre, luego el aluminio y finalmente el vidrio borosilicato) en el soporte horizontal del aparato de dilatación térmica. Un extremo del tubo se fijó rígidamente mediante la pinza de ajuste del soporte, mientras que el extremo libre se apoyó con cuidado sobre la aguja cilíndrica de rodamiento.
    \item \textbf{Alineación del dial indicador:} Se acopló la aguja indicadora al eje de la aguja de rodamiento y se verificó que el indicador estuviera orientado verticalmente hacia abajo, marcando la referencia de cero en el dial graduado angular del transportador de ángulos.
    \item \textbf{Medición de la longitud libre inicial:} Empleando la regla métrica graduada en milímetros, se midió la longitud inicial libre del tubo ($L_0$) desde el punto de sujeción fijo (pinza) hasta el punto de apoyo móvil sobre la aguja cilíndrica.
    \item \textbf{Registro de la temperatura ambiente:} Se midió la temperatura inicial del aire dentro del laboratorio ($T_0$) utilizando el termómetro de precisión, asumiendo que el tubo inicialmente se encontraba en equilibrio térmico con el entorno.
    \item \textbf{Preparación del generador de vapor:} Se llenó parcialmente el matraz de ebullición del generador de vapor con agua destilada y se encendió el elemento calefactor eléctrico. Durante esta fase de calentamiento inicial, la manguera flexible de salida de vapor se mantuvo desconectada del tubo de ensayo para evitar la pre-dilatación de la muestra.
    \item \textbf{Inyección de vapor de agua:} Una vez que el agua en el matraz alcanzó el punto de ebullición constante a 100$^\circ$C, se acopló la manguera de salida al conector de entrada del tubo de ensayo, permitiendo el flujo continuo de vapor caliente a través de su canal interior.
    \item \textbf{Monitoreo de la expansión térmica:} Conforme el vapor calentaba el tubo, se observó visualmente el desplazamiento gradual de la aguja indicadora sobre la escala angular del transportador. El paso del vapor se mantuvo de forma ininterrumpida hasta que el indicador angular dejó de moverse, señal de que la muestra había alcanzado el equilibrio térmico con el vapor a la temperatura final de $T_f = 100^\circ$C.
    \item \textbf{Registro de la variación angular:} Se leyó y registró con el transportador de ángulos el ángulo de giro final de la aguja ($\Delta\theta$) que mide el desplazamiento lineal acumulado.
    \item \textbf{Enfriamiento y cambio de muestra:} Se retiró la manguera de vapor, se dejó enfriar la estructura del dilatómetro a temperatura ambiente y se procedió a retirar la muestra ensayada. Este flujo de trabajo se repitió de manera independiente para cada uno de los materiales restantes.
\end{enumerate}